\scanpage{40}
Inversement, partant d'un groupe abélien $\Gamma^+$,
et d'un torseur $\Gamma^-$ sous $\Gamma^+$, on définit sur
$\Gamma=\Gamma^+\amalg\Gamma^-$ une \deleted{structure} loi de composition
par la formule précédente, et on voit de
suite que celle-ci est une loi de groupe, qui
fait de $\Gamma$ un groupe admettant $\Gamma^+$ comme
sous-groupe \inserted{abélien} d'indice 2, l'extension de $\mathbb Z_2$ par
$\Gamma^+$ ainsi obtenue donnant lieu à \deleted{l'op. triviale} \inserted{l'involution $x\mapsto-x$}
de $\Gamma^+$, et étant scindée par n'importe quel élément
de $\Gamma^-$. [Donc $\Gamma,\Gamma^+$ satisfont les conditions a) à d)
ci-dessus, mais pas néc. e)] On trouve
donc un foncteur qui va de la cat. des
\deleted{\unsure{gp}} torseurs sous groupes abéliens (avec les
morphismes habituels, pas néc iso) vers la catégorie
des ext. ``\unsure{dièdrales|diédrales}'' de $\mathbb Z/2$ par gpes abéliens
-- et c'est une équiv. de catégories -- et la
restriction de ce foncteur \deleted{aux} \inserted{à la ss-cat. \unsure{plénière|pleine} des} couples
$(\Gamma^+,\Gamma^-)$, avec $\Gamma^+$ satisfaisant e) ci-dessus, \inserted{et les seuls iso} est un
foncteur pl. fidèle vers la cat. des
groupes $\Gamma$ satisfaisant les conditions du début,
avec les iso. En particulier

\uline{Prop} Soit $n\geq3$ un entier. Le \deleted{cat.}\inserted{groupoïde} des formes de
$\mathbb D_n$ est équivalente au \deleted{la catégorie}\inserted{groupoïde} des droites
affines sur l'anneau \unsure{$k_n$|$\mathbb Z/n$}. En particulier, on
trouve

\[
\mathrm{Aut}(\mathbb D_n)\simeq\mathrm{Aff}(1,\mathbb Z/n)\simeq(\mathbb Z/n)^*\mathbin{\cdot_{1/2}}\mathbb Z/n
\]

% archive sheet number: 39
